\section{La ventana de contexto humana: el ser humano reduce la entropía, la IA la aumenta}

Esta sección parte de la «conversación de texto largo» con la que abre el episodio. Zhang Xiaojun (张小珺) le pregunta a Li Xiang (李想): si fueras un modelo grande, ¿qué tan grande sería tu ventana de contexto? La respuesta de Li Xiang no sigue la vía de las cifras de tokens, sino que se desplaza hacia la diferencia entre el ser humano y la IA: las personas no son buenas para procesar información muy compleja, así que deben reducir la entropía; la IA es buena para procesar volúmenes enormes de información, lo que se parece más a aumentar la entropía. Aquí, «reducir la entropía» no es la definición estricta de la física, sino un método de gestión y de cognición: comprimir un mundo complejo en unos cuantos juicios clave, unas cuantas acciones y unos cuantos principios.

\begin{figure}[H]
\centering
\includegraphics[width=0.94\textwidth]{human-ai-context-window.png}
\caption{Contexto humano frente a contexto de la IA: al ser humano le conviene más restar; a la IA, procesar información a gran escala. Diagrama conceptual propio, elaborado a partir de la conversación de 00:02:35--00:04:10.}
\end{figure}

\begin{knowledgebox}{Lectura de la figura: una ventana de contexto más grande no hace a nadie más parecido a un CEO}
Un modelo grande puede procesar muchos tokens, pero el valor de un CEO suele estar en la compresión: convertir información compleja en estrategia, organización y acción. A la inversa, si la IA se queda en «saber mucho» y no se conecta con herramientas, acciones y retroalimentación, tampoco puede generar valor real.
\end{knowledgebox}

\subsection{Tres niveles de herramientas: herramientas de información, herramientas de apoyo y herramientas de producción}

La sección anterior distinguió la capacidad de contexto del ser humano y la de la IA; esta sección da un paso más y pregunta cómo se convierte esa diferencia de capacidades en valor de producto. Li Xiang divide las herramientas en tres niveles. Las herramientas de información permiten a las personas encontrar información; las herramientas de apoyo aumentan su eficiencia; las herramientas de producción generan directamente una productividad por la que se paga. Esta clasificación es importante porque ofrece un criterio más firme para valorar los productos de IA: no se trata de si pueden responder, sino de si el usuario está dispuesto a pagar por ellos y a confiarles trabajo real.

\begin{importantbox}{Criterio para identificar una herramienta de producción}
Si una herramienta de IA solo hace que las personas «se sientan inteligentes», puede seguir siendo apenas una herramienta de información; solo cuando puede sustituir o liberar el trabajo real y frecuente de las personas, y el usuario está dispuesto a pagar por el resultado, empieza a entrar en el nivel de herramienta de producción.
\end{importantbox}

\subsection{DeepSeek: aplicar las mejores prácticas de la forma más sencilla}

A continuación, la entrevista se dirige a DeepSeek. Lo que Li Xiang aprendió de DeepSeek no es una técnica puntual, sino «aplicar de la forma más sencilla las mejores prácticas de la humanidad»: primero investigar, luego desarrollar y después convertirlo en negocio; si la brecha es de resultados, se ajusta la estrategia; si la brecha es de capacidades, se vuelve a la construcción de capacidades. Esta ruta parece sencilla, pero va contra la naturaleza humana, porque exige que la organización reconozca una y otra vez que sus capacidades son insuficientes, en lugar de avanzar a la fuerza solo con acciones de negocio.

\begin{figure}[H]
\centering
\includegraphics[width=0.96\textwidth]{deepseek-best-practice.png}
\caption{Las mejores prácticas aprendidas de DeepSeek: usar de la forma más sencilla las mejores prácticas de la humanidad y luego consolidar la organización y la ingeniería. Diagrama conceptual propio, elaborado a partir de la conversación de 00:15:01--00:19:10.}
\end{figure}

\begin{knowledgebox}{Lectura de la figura: DeepSeek como caso de organización y no solo como caso de modelo}
La secuencia de la figura, que va de MoE a la práctica, expresa una ruta de aprendizaje organizacional: primero investigar el problema, luego construir capacidades, después validarlas con ingeniería y, por último, convertir esas capacidades en resultados de negocio. A Li Xiang le interesa DeepSeek porque muestra cómo un equipo pequeño, una alta densidad de talento, el código abierto y la práctica de ingeniería pueden sumar fuerzas.
\end{knowledgebox}

\begin{warningbox}{Las «mejores prácticas» no producen automáticamente los mejores resultados}
Las mejores prácticas suelen ir contra la naturaleza humana, porque exigen bajar el ritmo para investigar, construir capacidades y hacer revisiones retrospectivas. Actuar según el impulso satisface más a la naturaleza humana, pero facilita confundir el problema con uno de estrategia de corto plazo y pasar por alto la verdadera brecha de capacidades.
\end{warningbox}

\subsection{Por qué desarrollar de todos modos capacidades propias de modelo base}

En la entrevista hay un detalle: dentro de Li Auto (理想) también existía la preocupación de que «usar modelos de código abierto podría perjudicar al equipo de desarrollo propio». La respuesta de Li Xiang es que un modelo de lenguaje de código abierto puede servir como base, pero lo que Li Auto quiere es desarrollar, a partir de sus escenarios de negocio, su propio VLA, su Agent OS y sus capacidades en el mundo físico. Lo central aquí no es elegir entre «código abierto o desarrollo propio», sino cómo se integra la base de código abierto en los datos, las tareas, la cadena de herramientas y los escenarios propios de la empresa.

\begin{knowledgebox}{Relación entre los modelos de código abierto y las capacidades de la empresa}
Linux puede ser de código abierto, y aun así Android formó sus propios componentes, API, ecosistema y experiencia de producto. Del mismo modo, los modelos de lenguaje pueden ser de código abierto, pero los fabricantes de automóviles siguen necesitando construir sus propias capacidades de sistema en torno al automóvil, los robots, los datos y los escenarios de los usuarios.
\end{knowledgebox}

\subsection{Resumen de la sección}

El capítulo inicial marca el tono de todo el episodio: la IA no es solo un contexto amplio y muchos parámetros, sino la manera de convertir la capacidad de procesar información en capacidad de actuar. DeepSeek se discute porque comprimió la investigación, la ingeniería y la organización en una metodología que se puede aprender.
